\section{Setting, investment and aggregation rules}
\label{sec:method}

\paragraph{Leaderboard and action.}
A leaderboard is a score matrix $X\in[0,100]^{M\times n}$: $M$ models, $n$ sub-datasets (``tasks''), higher is better. One model $\target$ is the \emph{target}; all other rows are held fixed (a single actor facing a static field). Targeted investment lets the target raise its score on a chosen task set $S\subseteq[n]$ to any $y_t\in[x_t,c_t]$ with $x=X_{\target,\cdot}$, and leave the other tasks unchanged; $c_t\ge x_t$ is a \emph{cap}. This models ``putting data relevant to sub-dataset $t$ into training'' as a score gain on $t$ only. It abstracts from the cost of obtaining the gain, from transfer between tasks, and from rivals' reactions; these assumptions and their limits are discussed in Section~\ref{sec:limits}.

\paragraph{First place.}
An aggregation rule $R$ maps the matrix to a ranking. The target is \emph{first} under $R$ if it is strictly ahead of every other model, with a margin of $\varepsilon=10^{-3}$ score points on the rule's own scale where the rule is cardinal (ties for first do not count). For rank-based rules ties in a task count half.

\paragraph{Investment: three coordinates.}
For a plan $y$ we measure the number of tasks touched $K=|\{t:y_t>x_t\}|$, the total score added $\Lone=\sum_t(y_t-x_t)$, and the largest single jump $\Linf=\max_t(y_t-x_t)$; $\Linf$ is the cost of a uniform improvement of the target on every task, the natural ``broad'' baseline against a targeted one. For each (target, rule, cap) we compute the plan that minimises $K$ and, separately, the plan that minimises $\Lone$ (and the smallest uniform raise for $\Linf$). The minimisers differ in general, so we always say which objective is minimised. If no plan within the caps makes the target first, the cell is \emph{infeasible}; infeasible cells are reported as a share and are never counted as zero or as maximal cost.

\paragraph{Caps (the main modelling assumption).}
How far a score can be pushed is not observable, so we treat it as a scenario: \texttt{ceil} ($c_t=100$, an upper envelope that is the most permissive), \texttt{front0} and \texttt{front2} (the best score of any other model on $t$, plus $0$ or $2$ points: nothing beyond the current frontier), \texttt{gain5}, \texttt{gain10}, \texttt{gain20} (target's own score plus $5/10/20$ points), and, for the Open LLM Leaderboard only, \texttt{emp95} and \texttt{emp99}, the 95th and 99th percentile per-task gains of fine-tuned models over their base models in the same leaderboard. Every result is shown per cap; none of the caps is claimed to be the truth.

\paragraph{Aggregation rules.}
\emph{Cardinal rules} aggregate the score vector of each model: mean (the rule of the Open LLM Leaderboard and the usual headline), geometric and harmonic mean, power means with exponent $p\in\{0.5,-0.5,-2,-5\}$, trimmed mean, median, and minimum (worst task). \emph{Ordinal rules} use only within-task comparisons: Borda count (equivalent to the mean rank), plurality (number of tasks with the highest score), the requirement of being the Condorcet winner (beating every rival in a strict majority of tasks), and Copeland. Being the Condorcet winner is sufficient for first place under every Condorcet-consistent rule (Copeland, Schulze, Kemeny, ranked pairs, minimax) but not necessary for any particular one, so its minimum cost is an \emph{upper bound} on the cost for such rules and is exact only for ``first under all of them''. We do not compute Kemeny, Schulze or ranked pairs.

\paragraph{Exact solvers.}
All minima are computed exactly, not by heuristics. For the mean the minimum $\Lone$ has the closed form $n(\tau+\varepsilon)-\sum_t x_t$ when caps do not bind, where $\tau$ is the best rival's mean, and the minimum $K$ is a greedy cover of the gap by the largest headrooms $c_t-x_t$. For power means with $p<1$ the optimum equalises $(y_t+s)$ at a water level; for the minimum and the median, $K$ is the number of tasks that must be lifted above the best rival's value (every task currently below it for the minimum; enough tasks to put a majority above it for the median, at most $\lfloor n/2\rfloor+1$). Trimmed mean and the ordinal rules are solved as mixed-integer programmes with lazily generated rival constraints, Borda, the Condorcet winner and plurality are validated against exhaustive enumeration on small subsamples of the real data (Table~\ref{tab:validation}); Copeland and the trimmed mean against brute force and a grid search in the test suite (\texttt{code/tests}). A solver failure or time limit is recorded as a failure and never as infeasible (none occurred in the HELM grid).
